\documentclass[quiet]{ctexart}

\ExplSyntaxOn \makeatletter

\keys_define:nn { hduthesis / exam }
  {
    watermark .tl_set:N = \l__hduthesis_exam_watermark_tl,
  }
\NewDocumentCommand \hdutexset {m} { \keys_set:nn { hduthesis / exam } {#1} }
\DeclareHookRule { shipout / background }
  { hduthesis / exam } { before } { pgfrcs }
\AddToHook { shipout / background } [ hduthesis / exam ]
  {
    \tl_if_empty:NF \l__hduthesis_exam_watermark_tl
      {
        \tikz [ remember~picture, overlay ]
          \node [ opacity = .2 ] at (current~page)
            {
              \includegraphics [ width = \linewidth ]
                { \l__hduthesis_exam_watermark_tl }
            };
      }
  }

\RequirePackage
  {
    geometry, mathtools, array, bm, esint, keytheorems, libertine, inconsolata,
    siunitx, fixdif, derivative, physics2, tasks, anyfontsize
  }
\RequirePackage[inline]{enumitem}
\RequirePackage{fontawesome5, fancyqr, circuitikz}
\hook_gput_code:nnn { cmd/faIcon/before } { . }
  {
    \group_begin:
    \fontsize { \fp_eval:n { .9 * \f@size } } { 0 } \selectfont
  }
\hook_gput_code:nnn { cmd/faIcon/after } { . } { \group_end: }
\geometry{margin = 1in}

\RequirePackage[svgnames]{xcolor}
\newkeytheorem{note}
  [ style = remark, numbered = false,
    name = \color{DarkRed}\sffamily \faIcon{info-circle}~Tip. ]
\newkeytheorem{problem}
  [ style = definition,
    name = \color{DarkGreen}\sffamily \faIcon{pen-square}~Problem ]
\newkeytheorem{solution}
  [ style = definition, numbered = false,
    name = \color{MidnightBlue}\sffamily \faIcon{check-square}~Solution ]
\usephysicsmodule{ab, op.legacy}
\usephysicsmodule{bm-um.legacy}
\settasks { label = \sffamily\Alph*., item-indent = 2.5em, label-offset = .5pt }
\newcommand \true {\addtocounter{task}{1}\faIcon{check-circle}}
\def \@result#1 {\underline{~#1~}}
\def \s@result {\textsf{【\qquad 】}}
\DeclareRobustCommand \result {\@ifstar\s@result\@result}
\keys_define:nn { textfig }
  {
    ratio      .tl_set:N    = \l__textfig_ratio_tl,
      ratio    .initial:n   = 7 \c_colon_str 2,
    position   .tl_set:N    = \l__textfig_position_tl,
      position .initial:n   = c,
    unknown    .code:n      = \tl_if_novalue:nF {#1}
      { \tl_set_eq:NN \l__textfig_position_tl \l_keys_key_tl }
  }
\NewDocumentEnvironment{textfig}{ O{} }
  {
    \group_begin:
    \keys_set:nn { textfig } {#1}
    \seq_set_split:Nee \l__whats_textfig_ratio_seq { \c_colon_str }
      { \l__textfig_ratio_tl }
    \noindent
    \minipage [ \l__textfig_position_tl ]
      {
        \fp_eval:n
          {
            .96*\seq_item:Nn \l__whats_textfig_ratio_seq {1}/
              (
                \seq_item:Nn \l__whats_textfig_ratio_seq {1} + 
                \seq_item:Nn \l__whats_textfig_ratio_seq {2}
              )
          }\dimexpr \linewidth
      }
  }
  {
    \endminipage
    \vspace{4.5pt plus 1.5pt minus 1.5pt} \par
    \group_end:
  }
\NewDocumentCommand \textfigsplit {}
  {
    \endminipage
    \hspace{.02\linewidth}
    \minipage [ \l__textfig_position_tl ]
      {
        \fp_eval:n
          {
            .96*\seq_item:Nn \l__whats_textfig_ratio_seq {2}/
              (
                \seq_item:Nn \l__whats_textfig_ratio_seq {1} + 
                \seq_item:Nn \l__whats_textfig_ratio_seq {2}
              )
          }\dimexpr \linewidth
      }
  }

% \hook_gput_code:nnn { cmd/maketitle/before } { . } { \vskip -\baselineskip }
% \hook_gput_code:nnn { cmd/maketitle/after } { . }
%   { \__weyoung_maketitle_aftercode_aux: }
\cs_new_protected:Nn \__weyoung_maketitle_aftercode_aux:
  {
    \vspace*{-2\baselineskip}
    \begin{center}
      \begin{minipage}{.66\textwidth}
        \centering
        \begin{tasks}
          [ label-width = 5\ccwd, item-indent = 2cm,
            after-item-skip = 0pt ](2)
          \task[\textbf{课程信息}：] 理学院/A0715022
          \task[\textbf{授课教师}：] 陈许敏
          \task[\textbf{解析提供}：] \textsf{Skyrmion}
          \task[{\makebox[5em][s]{\textbf{\textsf\TeX~排版}：}}] \textsf{Skyrmion}
        \end{tasks}
      \end{minipage}
      \begin{minipage}{.15\textwidth}
        \centering
        \begin{tikzpicture}
          \node [ inner~sep = 0 ] (wechat) {
            \fancyqr
              [ size = 4em, image~padding = 1,
                image = \raisebox{4em/33}{\textcolor{DarkGreen}{\faIcon{weixin}}},
                r~color = DarkGreen, version = 4
              ]{http://weixin.qq.com/r/sRJiemTEGjllrSAZ90ej}
            };
          \node [ inner~sep = .5ex, align = center, right, font = \scriptsize ]
            at (wechat.east) {未\\央\\学\\社};
        \end{tikzpicture}
      \end{minipage}
      \begin{minipage}{.15\textwidth}
        \centering
        \begin{tikzpicture}
          \node [ inner~sep = 0 ] (qq) {
            \fancyqr
              [ size = 4em, image~x~padding = 2, image~y~padding = 1,
                image = \raisebox{4em/33}{\textcolor{MidnightBlue}{\faIcon{qq}}},
                l~color = MidnightBlue, version = 4
              ]{https://qm.qq.com/q/PH7CsQWvm2}
            };
          \node [ inner~sep = .5ex, align = center, left, font = \scriptsize ]
            at (qq.west) {物\\理\\营};
        \end{tikzpicture}
      \end{minipage}
    \end{center}
  }

\newcommand \iu {{\mathup i\mkern1mu}}
\newcommand \eu {\mathup e}
\AtBeginDocument{%
  \let \Re \relax \let \Im \relax
  \DeclareMathOperator{\Re}{Re}
  \DeclareMathOperator{\Im}{Im}
}

\makeatother \ExplSyntaxOff
\usetikzlibrary{ fadings, patterns, arrows, decorations.markings }
\tikzset
  {
    > = stealth',
    every picture/.append style =
      { line join = round, line cap = round, thick }
  }

\setCJKmainfont{LXGW WenKai}[AutoFakeSlant]
\setCJKsansfont{LXGW Marker Gothic}[AutoFakeBold, AutoFakeSlant, Scale = .95]
\setCJKmonofont{LXGW WenKai Mono}[AutoFakeBold, AutoFakeSlant]
\ctexset{section/format+ = \raggedright}
\usepackage[warnings-off={mathtools-colon,mathtools-overbracket}]{unicode-math}
\setmathfont[AutoFakeBold]{STIX Two Math}

\makeatletter \def \@result #1{\underline{ \quad }} \makeatother
\DeclareMathOperator{\DFT}{DFT}

\title{\sffamily\bfseries 25Spring \textsc{Hdu} 「高等数学 A2」期中模拟}
\author{} \date{}

\begin{document}

\section{选择题}

\begin{problem}
  母线平行于$y$轴且通过曲线$\left\{ \begin{array}{r}
    x^{2} + 3y^{2} + z^{2} = 24 \\
    x^{2} - y^{2} + z^{2} = 0
    \end{array} \right.\ $的柱面方程为（）
    
    A.$4x^{2} + 3z^{2} = 24$ \quad B.$x^{2} + 3y^{2} = 24$ \quad C.$4x^{2} - z^{2} = 24$ \quad D.$4y^{2} - z^{2} = 24$
\end{problem}

\begin{problem}
  二元函数$z = f(x,y)$在$(x_{0},y_{0})$处可微是偏导数$f_{x}(x_{0},y_{0})$和$f_{y}(x_{0},y_{0})$存在的（）
        
        A.充分不必要条件 \quad B.必要不充分条件 \quad C.充要条件 \quad D.既不充分也不必要条件
\end{problem}

\begin{problem}
  设有直线$L_{1}: \frac{x - 2}{2} = \frac{y + 3}{- 1} = \frac{z - 1}{3}$，与$L_{2}: \left\{ \begin{array}{r}
    x + y = 4 \\
    y - z = 2
    \end{array} \right.\ $，则$L_{1}$与$L_{2}$的夹角为（）
    
    A.$\frac{\pi}{6}$ \quad B.$\frac{\pi}{4}$ \quad C.$\frac{\pi}{3}$ \quad D.$\frac{\pi}{2}$
\end{problem}

\begin{problem}
  函数$u = 2x^{2} + 3y^{2} + 4z^{2}$在点$(1,1,1)$处的方向导数最大值为（）
        
        A.$2\sqrt{11}$ \quad B.$2\sqrt{13}$ \quad C.$2\sqrt{29}$ \quad D.$2\sqrt{57}$
\end{problem}

\begin{problem}
  曲面$x^{2} + y^{2} + \cos{xy} - yz + xy = 0$上点$(0,1,2)$处的切平面方程（）
        
        A.$x + y - z = - 2$ \quad B.$x + z = 2$ \quad C.$x + 2y - z = 3$ \quad D.$y - 2z = 3$
\end{problem}

\begin{problem}
  函数$f(x,y) = 2x + (y - 2)\arccos\sqrt{\frac{x}{y}}$，则$f_{x}(x,2) =$（）
        
        A.$1$ \quad B.$2$ \quad C.$3$ \quad D.$4$
\end{problem}

\begin{problem}
  $\int_{0}^{2}{dx}\int_{x}^{2x}{f(x,y)dy} = \int_{0}^{\Psi}{dy}\int_{\varphi(y)}^{\Psi(y)}{f(x,y)dx}$成立，则$\varphi(y)$和$\Psi(y)$分别为（）
        
        A.$\varphi(y) = \frac{y}{2},\Psi(y) = y$ \quad B.$\varphi(y) = y,\Psi(y) = \frac{y}{2}$ \\
        C.$\varphi(y) = \frac{y}{2},\Psi(y) = 2y$ \quad D.$\varphi(y) = 2y,\Psi(y) = \frac{y}{2}$
\end{problem}

\begin{problem}
  曲线$\left\{ \begin{array}{r}
    x = 2t \\
    y = t^{2} \\
    z = - t^{3}
    \end{array} \right.\ $，在$t = 2$处的切线方程为（）
    
    A.$\frac{x - 4}{2} = \frac{y - 4}{4} = \frac{z + 8}{- 12}$ \quad B.$\frac{x - 4}{2} = \frac{y - 1}{1} = \frac{z + 1}{- 3}$ \\
    C.$(x - 4) + (y - 4) - (z - 8) = 0$ \quad D.$(x - 2) + (y - 1) - (z + 1) = 0$
\end{problem}

\begin{problem}
  设$f(x,y)$是连续函数，则$\int_{0}^{\frac{\pi}{2}}{d\theta}\int_{0}^{2\cos\theta}{f(\rho\cos\theta,\rho\sin\theta)\rho d\rho} =$（）
        
        A.$\int_{0}^{2}{dx}\int_{0}^{\sqrt{4 - x^{2}}}{f(x,y)dy}$ \quad B.$\int_{0}^{2}{dy}\int_{0}^{\sqrt{4 - y^{2}}}{f(x,y)dx}$ \\
        C.$\int_{0}^{1}{dx}\int_{0}^{\sqrt{2x - x^{2}}}{f(x,y)dy}$ \quad D.$\int_{0}^{1}{dy}\int_{0}^{\sqrt{2y - y^{2}}}{f(x,y)dx}$
\end{problem}

\begin{problem}
  已知函数$f(x,y)$在点$(x_{0},y_{0})$的某邻域内满足有一阶和二阶连续的偏导数，且$f_{x}(x_{0},y_{0}) = f_{y}(x_{0},y_{0}) = 0$，记$A = f_{xx}(x_{0},y_{0})$，$B = f_{xy}(x_{0},y_{0})$，$C = f_{yy}(x_{0},y_{0})$，则满足（）时函数在点$(x_{0},y_{0})$处取得极小值
        
        A.$AC - B^{2} < 0$且$A > 0$ \quad B.$AC - B^{2} < 0$且$A < 0$ \\
        C.$AC - B^{2} > 0$且$A > 0$ \quad D.$AC - B^{2} > 0$且$A < 0$
\end{problem}

\section{填空题}

\begin{problem}
  经过已知点$(1,-1,4)$和直线$\frac{2}{x+1}=\frac{y}{5}=\frac{1-z}{-1}$的平面方程是（\rule{3cm}{0.15mm}）
\end{problem}

\begin{problem}
  曲线$\{3x^{2} + 2y^{2} = 12,z=0\}$绕$y$轴旋转一周达到的旋转面在点$(0,\sqrt{3},\sqrt{2})$处的指向外侧的单位法向量为（\rule{3cm}{0.15mm}）
\end{problem}

\begin{problem}
  设$f(x,y) = \int_{0}^{xy}e^{xt^{2}}dt$，则$\left.\frac{\partial^2 f}{\partial x \partial y}\right|_{(1,1)}=$（\rule{3cm}{0.15mm}）
\end{problem}

\begin{problem}
  
\end{problem}

\begin{enumerate}

    \item[14.] 设曲线$u(x,y,z)=1+\frac{x^2}{6}+\frac{y^2}{12}+\frac{z^2}{18}$，单位向量$\mathbf{n}=\frac{1}{\sqrt{3}}(1,1,1)$，则$\left.\frac{\partial u}{\partial n}\right|_{(1,2,3)}$等于（\rule{3cm}{0.15mm}）
\end{enumerate}

三、计算题

\begin{enumerate}
    \item[15.] 求过点$(-1,0,4)$，且平行于平面$3x - 4y + z - 10 = 0$，又与直线$\frac{x + 1}{1} = \frac{y - 3}{1} = \frac{z}{2}$相交的直线的方程。

    \item[16.] 设平面$\pi$：$x + y + 2z - 1 = 0$与球面$x^{2} + y^{2} + z^{2} = 1$相交成一个圆，试求这个圆的圆心。

    \item[17.] 求过点$A(2, -3,1)$且与直线$\frac{x - 1}{2} = \frac{y + 1}{-1} = \frac{z - 3}{3}$垂直相交的直线方程。

    \item[18.] 求平面$\frac{x}{3} + \frac{y}{4} + \frac{z}{5} = 1$和柱面$x^{2} + y^{2} = 1$的交线上与$xOy$面距离最短的点。

    \item[19.] 某厂家生产的一种产品同时在两个市场销售，售价分别为$p_1$和$p_2$，销售量分别为$q_1$和$q_2$，需求函数分别为$q_1 = 24 - 0.2p_1$，$q_2 = 10 - 0.05p_2$，总成本函数为$C = 35 + 40(q_1 + q_2)$。试问：厂家如何确定两个市场的售价，能使其获得的总利润最大，最大总利润为多少？

    \item[20.] 求函数$f(x,y) = x^{3} + 8y^{3} - xy$的极值。

    \item[21.] 计算以$xOy$面上的圆周$x^{2} + y^{2} = ax\ (a > 0)$围成的闭区域为底，而以曲面$z = x^{2} + y^{2}$为顶的曲顶柱体的体积。

    \item[22.] 计算三重积分$\iiint_{\Omega}\sqrt{y^{2} + z^{2}}\,dxdydz$，其中$\Omega$由$xOy$平面上的曲线$x = hy\ (h > 0)$绕$x$轴旋转而成的曲面与平面$x = h$所围成的闭区域。

    \item[23.] 设函数$f(x)$在$[a,b]$上连续，证明不等式：
    \[
    \left( \int_{a}^{b}f(x)\,dx \right)^{2} \leq (b - a)\int_{a}^{b}f^{2}(x)\,dx
    \]
\end{enumerate}

\end{document}